\input{../tex-inputs/latex-leseansicht-vorspann}

\section[Arthur Schnitzler an Gustav Schwarzkopf, 9. 7. 1915]{L04489 Arthur Schnitzler an Gustav Schwarzkopf, 9. 7. 1915}
\nopagebreak\mylabel{L04489v}
\rehead{ }\normalsize\beginnumbering\briefempfaengerindex{Schwarzkopf, Gustav@\textsc{Schwarzkopf, Gustav}!zzzSchnitzler, Arthur@\emph{von Arthur Schnitzler}!1915-07-091@{9. 7. 1915}|(be}
\toendnotes[C]{\smallbreak\pagebreak[2]}
\correspDesc{Versand  durch Arthur Schnitzler am 9. 7. 1915 in Wien
\newline{}Übermittlung  am 9. 7. 1915 in Bad Ischl
\newline{}Erhalt  durch Gustav Schwarzkopf im Zeitraum [10. 7. 1915 – 14. 7. 1915?] in Wien}\toendnotes[C]{\smallbreak}
\Standort{DLA, A:Schnitzler, HS.1985.1.1897.}
\physDesc{Bildpostkarte, 674 Zeichen
\newline{}Handschrift: schwarze Tinte, deutsche Kurrent
\newline{}Versand: Stempel: »\nobreak{}\oindex{Wien@\textbf{Wien}|pwk}18/1 Wien 110, 9. VII. 15, 3\nobreak{}«.  }\toendnotes[C]{\smallbreak}\pstart{}{\pb}Herrn Guſtav Schwarzkopf\pend{}\pstart{}Iſchl\oindex{Bad Ischl@\textbf{Bad Ischl}|pw}\pend{}\pstart{}\textsc{Hotel garni} Athen\oindex{Hotel garni Athen@\textbf{Hotel garni Athen}|pw}\pend{}{\bigskip}
\pstart
           \noindent{}{\pb}\textcolor{gray}{\textbf{Wien, XVIII, Sternwartestr. 71\oindex{Wien@\textbf{Wien}!XVIII., Währing@\textbf{XVIII., Währing}!Sternwartestraße 71@\textbf{Sternwartestraße 71}|pw}.}}\pend
           \vspace{1em}
\pstart
           \raggedleft{}9/7 1915\pend
           \vspace{0.5em}
\pstart
           lieber Guſtav, danke für \label{K_L04489-1v}\edtext{Nachricht u Gruſs}{\lemma{\textnormal{\emph{Nachricht u Gruss}}}\Cendnote{\textnormal{XXXX Auszeichnungsfehler: Dokument L04246 nicht gefunden.}}}\label{K_L04489-1}. Laſſen Sie{ }ſichs weiter
               wohlergehn u mich hoffen, daſs ich Sie am 28. od 29. noch
               dort finde. O\pwindex{Schnitzler, Olga 17.\,1.\,1882 Wien – 13.\,1.\,1970 Lugano@\textsc{Schnitzler, Olga} (17.\,1.\,1882 Wien – 13.\,1.\,1970 Lugano)|pw} fährt \label{K_L04489-2v}\edtext{am 17.}{\lemma{\textnormal{\emph{am 17.}}}\Cendnote{\textnormal{Laut \emph{Tagebuch}\pwindex{Schnitzler, Arthur 15. 5. 1862 Wien – 21. 10. 1931 ebd.@\textsc{Schnitzler, Arthur} (15. 5. 1862 Wien – 21. 10. 1931 ebd.)!Tagebuch@\strich\emph{Tagebuch}|pwk} brachen die beiden\pwindex{Schnitzler, Olga 17.\,1.\,1882 Wien – 13.\,1.\,1970 Lugano@\textsc{Schnitzler, Olga} (17.\,1.\,1882 Wien – 13.\,1.\,1970 Lugano)|pwkv}\pwindex{Bachrach, Stefanie 22.\,5.\,1887 Wien – 16.\,5.\,1917 ebd.@\textsc{Bachrach, Stefanie} (22.\,5.\,1887 Wien – 16.\,5.\,1917 ebd.)|pwkv} bereits am 15. 7. 1915 auf.}}}\label{K_L04489-2} mit Ste\textcolor{gray}{ph}i\pwindex{Bachrach, Stefanie 22.\,5.\,1887 Wien – 16.\,5.\,1917 ebd.@\textsc{Bachrach, Stefanie} (22.\,5.\,1887 Wien – 16.\,5.\,1917 ebd.)|pw} direct nach Altauſſee\oindex{Altaussee@\textbf{Altaussee}|pw} und gruppirt sich da{\geminationn}
               zu mir um. Vorgeſtern haben wir in Kritzendorf\oindex{Kritzendorf@\textbf{Kritzendorf}|pw}{ }ſtrand u ſtromgebadet\oindex{Strombad Kritzendorf@\textbf{Strombad Kritzendorf}|pwv}. Ein
               herrliches Gewitter mit Blitzen von denen Lili\pwindex{Cappellini, Lili 13.\,9.\,1909 Wien – 26.\,7.\,1928 Venedig@\textsc{Cappellini, Lili} (13.\,9.\,1909 Wien – 26.\,7.\,1928 Venedig)|pw}{ }ſagte: {[}»{]}Jetzt machen {\pb}die Engeln die
               Thür auf« brachte geſtern Abend nach einem der{ }ſchwülſten Tage die ich je erlebt,
                     \textcolor{gray}{willko{\geminationm}ne} Frische in Haus, 
               Garten u Welt. Ich find es überhaupt in Wien\oindex{Wien@\textbf{Wien}|pw}{ }ſehr{ }ſchön. Grüßen Sie Richard\pwindex{Beer-Hofmann, Richard 11.\,7.\,1866 Wien – 26.\,9.\,1945 New York City@\textsc{Beer-Hofmann, Richard} (11.\,7.\,1866 Wien – 26.\,9.\,1945 New York City)|pw} und die Seinen\pwindex{Beer-Hofmann, Paula 25.\,2.\,1879 Wien – 30.\,10.\,1939 Zürich@\textsc{Beer-Hofmann, Paula} (25.\,2.\,1879 Wien – 30.\,10.\,1939 Zürich)|pw}\pwindex{Beer-Hofmann, Gabriel 9.\,1.\,1901 Wien – 24.\,3.\,1971 St Albans@\textsc{Beer-Hofmann, Gabriel} (9.\,1.\,1901 Wien – 24.\,3.\,1971 St Albans)|pw}\pwindex{Beer-Hofmann, Naëmah 20.\,12.\,1898 Wien – 10.\,11.\,1971 New York City@\textsc{Beer-Hofmann, Naëmah} (20.\,12.\,1898 Wien – 10.\,11.\,1971 New York City)|pw}\pwindex{Beer-Hofmann, Mirjam 4.\,9.\,1897 Wien – 24.\,12.\,1984 New York City@\textsc{Beer-Hofmann, Mirjam} (4.\,9.\,1897 Wien – 24.\,12.\,1984 New York City)|pw}. Ich hoffe \textcolor{gray}{Jacobs}{ }\textcolor{gray}{leiter} beſteigen zu dürfen – höhere \textcolor{gray}{Touren} mach
               ich nicht.\pend
           
\pstart
           herzlichſt{\\[\baselineskip]}Ihr{\\[\baselineskip]}\spacefill\mbox{A.}\pend
           \leftskip=0em{}\selectlanguage{ngerman}\endnumbering\briefempfaengerindex{Schwarzkopf, Gustav@\textsc{Schwarzkopf, Gustav}!zzzSchnitzler, Arthur@\emph{von Arthur Schnitzler}!1915-07-091@{9. 7. 1915}|)be}\mylabel{L04489h}
\begin{anhang}
\end{anhang}\newcommand{\dateiname}{L04489}\newcommand{\titel}{Arthur Schnitzler an Gustav Schwarzkopf, 9. 7. 1915}\newcommand{\editorInnen}{Herausgegeben von Jahnke, SelmaMüller, Martin Anton}\input{../tex-inputs/latex-leseansicht-abspann}
